% ============================================================
% 07_CONCLUSAO.TEX
% Sistema Adaptativo de Classificação Multicritério
% de Fluidos de Corte
%
% FINALIDADE
%
% Estruturar a conclusão da apresentação.
%
% Esta seção deverá sintetizar:
%
% - resposta ao problema;
% - resultado principal;
% - robustez da conclusão;
% - limitações;
% - implicações práticas;
% - mensagem final.
%
% Este arquivo constitui apenas a BASE da apresentação.
%
% Nenhum vencedor, ranking, conclusão científica ou recomendação
% deve ser preenchido antes da execução completa das análises.
%
% ============================================================


% ============================================================
% FRAME 1 — RETOMADA DO OBJETIVO
% ============================================================

\begin{frame}{Retomando o objetivo}

    \begin{block}{Problema}
        Classificar os fluidos de corte considerando simultaneamente
        múltiplos critérios potencialmente conflitantes, sem reduzir
        a decisão a uma única medida.
    \end{block}

    \vspace{0.3cm}

    O projeto foi estruturado para produzir uma classificação:

    \begin{itemize}
        \item transparente;
        \item reproduzível;
        \item rastreável;
        \item multicritério;
        \item acompanhada de análise de robustez.
    \end{itemize}

    \vspace{0.3cm}

    \begin{block}{Estratégia}
        Auditoria dos dados
        $\rightarrow$
        critérios
        $\rightarrow$
        CRITIC
        $\rightarrow$
        TOPSIS
        $\rightarrow$
        Pareto
        $\rightarrow$
        sensibilidade
        $\rightarrow$
        robustez.
    \end{block}

\end{frame}


% ============================================================
% FRAME 2 — PRINCIPAL RESULTADO
% ============================================================

\begin{frame}{Principal resultado}

    \begin{block}{Classificação principal}
        \texttt{TO\_BE\_FILLED}
    \end{block}

    \vspace{0.3cm}

    \begin{block}{Primeira alternativa}
        \texttt{TO\_BE\_FILLED}
    \end{block}

    \vspace{0.3cm}

    \begin{block}{Diferença para as alternativas seguintes}
        \texttt{TO\_BE\_FILLED}
    \end{block}

    \vspace{0.3cm}

    \begin{alertblock}{Regra}
        Esta conclusão só deverá ser preenchida depois da aprovação
        do ranking oficial produzido por Vitor.
    \end{alertblock}

\end{frame}


% ============================================================
% FRAME 3 — CONCLUSÃO INTEGRADA
% ============================================================

\begin{frame}{Conclusão integrada}

    A recomendação final não deverá considerar somente a posição
    no ranking.

    \vspace{0.3cm}

    Ela deverá integrar:

    \begin{itemize}
        \item coeficiente TOPSIS;
        \item diferença entre as primeiras posições;
        \item conformidade técnica;
        \item comportamento nos cenários;
        \item posição na análise de Pareto;
        \item sensibilidade aos pesos;
        \item sensibilidade à normalização;
        \item estabilidade frente à retirada de critérios;
        \item estabilidade frente à retirada de alternativas.
    \end{itemize}

\end{frame}


% ============================================================
% FRAME 4 — ROBUSTEZ DA CONCLUSÃO
% ============================================================

\begin{frame}{Robustez da conclusão}

    \begin{block}{Robustez local}
        \texttt{TO\_BE\_FILLED}
    \end{block}

    \vspace{0.2cm}

    \begin{block}{Robustez ampliada}
        \texttt{TO\_BE\_FILLED}
    \end{block}

    \vspace{0.2cm}

    \begin{block}{Estabilidade do top 3}
        \texttt{TO\_BE\_FILLED}
    \end{block}

    \vspace{0.2cm}

    \begin{block}{Vencedor robusto}
        \texttt{TO\_BE\_FILLED}
    \end{block}

\end{frame}


% ============================================================
% FRAME 5 — POSSÍVEIS FORMAS DE CONCLUSÃO
% ============================================================

\begin{frame}{Como interpretar a conclusão?}

    A força da conclusão dependerá do comportamento observado nas
    análises de sensibilidade.

    \vspace{0.3cm}

    A conclusão poderá assumir diferentes formas.

    \begin{columns}[T]

        \begin{column}{0.48\textwidth}

            \begin{block}{Liderança estável}
                Uma alternativa mantém a primeira posição sob as
                principais configurações avaliadas.
            \end{block}

            \begin{block}{Top 3 estável}
                As mesmas alternativas permanecem entre as primeiras,
                embora sua ordem possa variar.
            \end{block}

        \end{column}


        \begin{column}{0.48\textwidth}

            \begin{block}{Liderança sensível}
                Alterações metodológicas razoáveis modificam o primeiro
                colocado.
            \end{block}

            \begin{block}{Sem vencedor único}
                As evidências sustentam um grupo de alternativas
                competitivas em vez de uma única superior.
            \end{block}

        \end{column}

    \end{columns}

    \vspace{0.2cm}

    \centering
    \textit{A interpretação correspondente será escolhida somente após
    a execução.}

\end{frame}


% ============================================================
% FRAME 6 — CONFORMIDADE NA CONCLUSÃO
% ============================================================

\begin{frame}{Desempenho e conformidade}

    \begin{block}{Pergunta final}
        A alternativa melhor classificada também atende às exigências
        técnicas confirmadas?
    \end{block}

    \vspace{0.4cm}

    Status do primeiro colocado:

    \begin{center}
        \texttt{TO\_BE\_FILLED}
    \end{center}

    \vspace{0.3cm}

    Restrições relevantes:

    \begin{center}
        \texttt{TO\_BE\_FILLED}
    \end{center}

    \vspace{0.3cm}

    \begin{alertblock}{Importante}
        O ranking multicritério não deverá ser utilizado para ocultar
        eventual não conformidade técnica.
    \end{alertblock}

\end{frame}


% ============================================================
% FRAME 7 — PARETO NA CONCLUSÃO
% ============================================================

\begin{frame}{Consistência com a análise de Pareto}

    \begin{block}{Status do primeiro colocado}
        \texttt{TO\_BE\_FILLED}
    \end{block}

    \vspace{0.3cm}

    \begin{block}{Alternativas não dominadas}
        \texttt{TO\_BE\_FILLED}
    \end{block}

    \vspace{0.3cm}

    \begin{block}{Interpretação}
        \texttt{TO\_BE\_FILLED}
    \end{block}

    \vspace{0.2cm}

    \begin{alertblock}{Cuidado}
        Pareto complementa a interpretação, mas não substitui o ranking
        multicritério.
    \end{alertblock}

\end{frame}


% ============================================================
% FRAME 8 — LIMITAÇÕES
% ============================================================

\begin{frame}{Limitações}

    \begin{block}{Limitações confirmadas}
        \begin{itemize}
            \item \texttt{TO\_BE\_FILLED}
            \item \texttt{TO\_BE\_FILLED}
            \item \texttt{TO\_BE\_FILLED}
        \end{itemize}
    \end{block}

    \vspace{0.3cm}

    As limitações apresentadas deverão ser apenas aquelas
    efetivamente identificadas durante o projeto.

    \vspace{0.2cm}

    \begin{alertblock}{Regra}
        Não inserir limitações genéricas apenas para completar o slide.
    \end{alertblock}

\end{frame}


% ============================================================
% FRAME 9 — IMPACTO DAS LIMITAÇÕES
% ============================================================

\begin{frame}{Impacto das limitações}

    Para cada limitação importante, a apresentação deverá responder:

    \begin{itemize}
        \item qual parte da análise é afetada?
        \item qual interpretação exige cautela?
        \item ela altera o ranking?
        \item ela altera a confiança na conclusão?
        \item ela exige coleta adicional de dados?
    \end{itemize}

    \vspace{0.3cm}

    \begin{block}{Síntese}
        \texttt{TO\_BE\_FILLED}
    \end{block}

\end{frame}


% ============================================================
% FRAME 10 — RECOMENDAÇÃO
% ============================================================

\begin{frame}{Recomendação}

    \begin{block}{Recomendação principal}
        \texttt{TO\_BE\_FILLED}
    \end{block}

    \vspace{0.3cm}

    \begin{block}{Condições ou ressalvas}
        \texttt{TO\_BE\_FILLED}
    \end{block}

    \vspace{0.3cm}

    \begin{alertblock}{Regra}
        A recomendação deverá ser proporcional à força das evidências
        obtidas nas análises de robustez.
    \end{alertblock}

\end{frame}


% ============================================================
% FRAME 11 — RECOMENDAÇÃO CONDICIONAL
% ============================================================

\begin{frame}{Recomendações condicionais}

    Caso os cenários mostrem diferenças relevantes, poderão ser
    apresentadas recomendações condicionais.

    \vspace{0.3cm}

    \begin{block}{Cenário global}
        \texttt{TO\_BE\_FILLED}
    \end{block}

    \begin{block}{Cenário alternativo 1}
        \texttt{TO\_BE\_FILLED}
    \end{block}

    \begin{block}{Cenário alternativo 2}
        \texttt{TO\_BE\_FILLED}
    \end{block}

    \vspace{0.2cm}

    \centering
    \textit{Manter este slide somente se os cenários realmente
    acrescentarem informação à decisão.}

\end{frame}


% ============================================================
% FRAME 12 — CONTRIBUIÇÕES DO PROJETO
% ============================================================

\begin{frame}{Contribuições do projeto}

    O trabalho pretende contribuir com uma estrutura de decisão que:

    \begin{itemize}
        \item começa pela qualidade dos dados;
        \item documenta as decisões de tratamento;
        \item considera simultaneamente múltiplos critérios;
        \item utiliza ponderação objetiva;
        \item produz uma classificação interpretável;
        \item separa desempenho de conformidade;
        \item verifica a estabilidade das conclusões;
        \item mantém rastreabilidade da análise.
    \end{itemize}

\end{frame}


% ============================================================
% FRAME 13 — REPRODUTIBILIDADE
% ============================================================

\begin{frame}{Reprodutibilidade}

    A estrutura do projeto foi planejada para permitir:

    \begin{itemize}
        \item reconstrução da base tratada;
        \item repetição das análises;
        \item rastreamento das decisões;
        \item validação dos resultados;
        \item atualização do ranking caso os dados mudem.
    \end{itemize}

    \vspace{0.3cm}

    \begin{block}{Princípio}
        Alterar uma entrada relevante deverá provocar a reexecução
        das etapas dependentes, e não uma atualização manual dos
        resultados.
    \end{block}

\end{frame}


% ============================================================
% FRAME 14 — MENSAGEM FINAL
% ============================================================

\begin{frame}{Mensagem final}

    \begin{center}

        \Large

        \textbf{A escolha de um fluido não deve depender apenas de
        qual alternativa ocupa a primeira posição.}

        \vspace{0.6cm}

        \normalsize

        A decisão deve considerar:

        \vspace{0.3cm}

        \textbf{
        qualidade dos dados
        +
        desempenho multicritério
        +
        conformidade
        +
        robustez
        }

    \end{center}

\end{frame}


% ============================================================
% FRAME 15 — CONCLUSÃO FINAL A SER PREENCHIDA
% ============================================================

\begin{frame}{Conclusão}

    \begin{block}{Resultado}
        \texttt{TO\_BE\_FILLED}
    \end{block}

    \begin{block}{Robustez}
        \texttt{TO\_BE\_FILLED}
    \end{block}

    \begin{block}{Limitação principal}
        \texttt{TO\_BE\_FILLED}
    \end{block}

    \begin{block}{Recomendação}
        \texttt{TO\_BE\_FILLED}
    \end{block}

\end{frame}


% ============================================================
% FRAME 16 — ENCERRAMENTO
% ============================================================

\begin{frame}[plain]

    \begin{center}

        \vfill

        {\LARGE\bfseries Obrigado!}

        \vspace{0.8cm}

        {\large Perguntas?}

        \vfill

    \end{center}

\end{frame}


% ============================================================
% REGRAS PARA O PREENCHIMENTO FUTURO
% ============================================================
%
% A conclusão final deve ser obtida a partir da cadeia:
%
% base aprovada
% ↓
% critérios aprovados
% ↓
% CRITIC validado
% ↓
% TOPSIS validado
% ↓
% cenários
% ↓
% Pareto
% ↓
% sensibilidade
% ↓
% robustez
%
% Não utilizar como conclusão final resultados de inspeções
% preliminares da planilha.
%
% ============================================================


% ============================================================
% NÃO FORÇAR UM VENCEDOR
% ============================================================
%
% Caso as análises mostrem instabilidade relevante, a conclusão
% poderá ser:
%
% - ausência de vencedor robusto;
% - grupo de melhores alternativas;
% - recomendação dependente do cenário.
%
% Essa também é uma conclusão cientificamente válida.
%
% ============================================================


% ============================================================
% COERÊNCIA COM O RELATÓRIO
% ============================================================
%
% Antes da apresentação:
%
% [ ] mesmo primeiro colocado
% [ ] mesmo top 3
% [ ] mesmos pesos
% [ ] mesma conformidade
% [ ] mesmo Pareto
% [ ] mesma robustez
% [ ] mesmas limitações
% [ ] mesma conclusão
%
% que no relatório técnico final.
%
% ============================================================


% ============================================================
% MATERIAL EXTERNO
% ============================================================
%
% Caso a professora forneça orientação específica para:
%
% - forma de conclusão;
% - conteúdo obrigatório;
% - terminologia;
% - recomendação;
%
% utilizar a orientação fornecida por ela.
%
% ============================================================


% ============================================================
% ESTADO ATUAL
% ============================================================
%
% CONCLUSION_SECTION_STATUS = BASE_PREPARED
%
% FINAL_RANKING = TO_BE_FILLED
% TOP_1 = TO_BE_FILLED
% ROBUSTNESS_CONCLUSION = TO_BE_FILLED
% FINAL_LIMITATIONS = TO_BE_FILLED
% FINAL_RECOMMENDATION = TO_BE_FILLED
%
% ============================================================